% Generated draft: requires mathematical and source review.
\documentclass[11pt]{article}
\usepackage[margin=25mm]{geometry}
\usepackage{fontspec}
\setmainfont{Latin Modern Roman}
\usepackage{amsmath,amssymb,mathtools,unicode-math}
\setmathfont{Latin Modern Math}
\usepackage{enumitem,needspace,xurl,hyperref}
\setlength{\parindent}{0pt}
\setlength{\parskip}{5pt}
\setlength{\emergencystretch}{3em}
\newcommand{\N}{\mathbb{N}}
\newcommand{\Z}{\mathbb{Z}}
\newcommand{\Q}{\mathbb{Q}}
\newcommand{\R}{\mathbb{R}}
\newcommand{\C}{\mathbb{C}}
\newcommand{\F}{\mathbb{F}}
\newcommand{\A}{\mathbb{A}}
\newcommand{\PP}{\mathbb{P}}
\newcommand{\E}{\mathbb{E}}
\newcommand{\eps}{\varepsilon}
\newcommand{\dd}{\,\mathrm{d}}
\begin{document}
\section*{UnsolvedMath ID: 701721}
\textbf{Source problem code:} catalogue entry
\subsection{Definitions and mathematical statement}
A group \(G\), interpreted here as discrete, is \emph{unitarisable} if, for every complex Hilbert space \(H\) and every representation \(\pi:G\to\operatorname{GL}(H)\) by bounded invertible operators satisfying
\[
\sup_{g\in G}\|\pi(g)\|<\infty,
\]
there is a bounded invertible operator \(S:H\to H\) such that \(S\pi(g)S^{-1}\) is unitary for every \(g\in G\). The source supplies no topology on \(G\).

For a countable discrete group \(G\), let \(G\curvearrowright(X,\mu)\) be an essentially free probability-measure-preserving action on a standard nonatomic probability space, and let
\[
\mathcal R_G=\{(x,gx):x\in X,\ g\in G\}
\]
be its orbit equivalence relation. A graphing of \(\mathcal R_G\) is a countable family \(\Phi=(\varphi_i:A_i\to B_i)_i\) of measure-preserving measurable partial bijections whose graphs generate \(\mathcal R_G\), modulo null sets. Define
\[
\operatorname{Cost}(\mathcal R_G)
=\inf_{\Phi}\sum_i\mu(A_i),
\qquad
\operatorname{Cost}(G)
=\inf_{G\curvearrowright(X,\mu)}\operatorname{Cost}(\mathcal R_G),
\]
where the second infimum ranges over such essentially free actions.

\textbf{(a)} The source asks whether every unitarisable group has \emph{trivial cost}. For countably infinite discrete groups, two possible precise readings are
\[
\forall G\quad
\bigl(G\text{ unitarisable}\Longrightarrow\operatorname{Cost}(G)=1\bigr)
\]
and the stronger fixed-price assertion
\[
\forall G\quad
\left(G\text{ unitarisable}\Longrightarrow
\left[\operatorname{Cost}(\mathcal R_G)=1
\text{ for every essentially free probability-measure-preserving action of }G\right]\right).
\]
In these two displays, \(G\) ranges over countably infinite discrete groups. The supplied record does not define ``trivial cost'' or specify which quantification over actions is intended; these are alternative readings, not an asserted equivalence. It also does not impose countability or infinitude. For finite groups the usual cost of a free action is \(1-1/|G|\), so a formulation covering finite groups must address this boundary case instead of requiring cost \(1\) universally. Thus part (a) retains an unresolved scope and terminology qualification.

Let \(\mathcal N(G)\) be the group von Neumann algebra of a discrete group \(G\), generated by its left regular representation on \(\ell^2(G)\), equipped with its canonical trace. With \(\C\) the trivial left \(\C G\)-module, define the first \(L^2\)-Betti number by
\[
\beta_1^{(2)}(G)
=\dim_{\mathcal N(G)}
\operatorname{Tor}^{\C G}_1\bigl(\mathcal N(G),\C\bigr),
\]
using the extended von Neumann dimension and the natural right \(\C G\)-module structure on \(\mathcal N(G)\).

\textbf{(b)} Under this discrete-group interpretation, the second requested assertion is
\[
\forall G\quad
\bigl(G\text{ unitarisable}\Longrightarrow\beta_1^{(2)}(G)=0\bigr).
\]
The source presents this as the weaker target to prove at least.

A group is residually finite if each nonidentity element has nonidentity image under some homomorphism to a finite group. The supplied source reports that part (a) is known for residually finite groups, citing I. Epstein and N. Monod. This report is preserved as a source qualification, without independently verifying it or resolving the intended meaning of ``trivial cost.''
\subsection{Short English statement}
The problem asks whether every unitarisable group has trivial cost and, at least, whether its first \(L^2\)-Betti number vanishes. The source reports the cost assertion for residually finite groups. The meaning of ``trivial cost,'' the quantification over actions, and the intended group-size restrictions require clarification.
\subsection{Sources}
\begin{itemize}
\item Local UnsolvedMath registry, supplied complete source record: permanent ID \(701721\), exact source problem code \texttt{catalogue entry}, title \emph{Kourovka Notebook Problem 20.63}. The supplied registry labels the problem \texttt{open}; this is a supplied status label, not an independently verified literature assessment.
\item \emph{The Kourovka Notebook (Unsolved Problems in Group Theory)}, 20th issue (2022), original problem 20.63, proposed by N. Monod. The supplied background also identifies the 21st issue edition (2026), \texttt{arXiv:1401.0300v48}. The supplied collection description names E. I. Khukhro and V. D. Mazurov as editors.
\item I. Epstein and N. Monod, \emph{Int. Math. Res. Notes}, \textbf{2009}, no. 22 (2009), 4336--4353. This is the citation supplied for the reported residually finite case of part (a). No article title or URL is supplied, and the citation and its precise scope have not been independently checked.
\end{itemize}
\end{document}
